\documentclass[book.tex]{subfiles}

\begin{document}

\chapter{Gauge Equivariant Neural Networks}
\label{chap:gauge_equivariant_networks}
\noindent\rule{11cm}{0.4pt} \\
\textbf{Prerequisites:}  \autoref{chap:ml_basics}, \autoref{chap:grad_descent}  \\
\textbf{Difficulty Level:} ** \\
\noindent\rule{11cm}{0.4pt}
\vspace{5mm}

%"""
%The principle of equivariance to symmetry transformations enables a theoretically grounded approach to neural network architecture design. Equivariant networks have shown excellent performance and data efficiency on vision and medical imaging problems that exhibit symmetries.

%Here we show how this principle can be extended beyond global symmetries to local gauge transformations. This enables the development of a very general class of convolutional neural networks on manifolds that depend only on the intrinsic geometry, and which includes many popular methods from equivariant and geometric deep learning. We implement gauge equivariant CNNs for signals defined on the surface of the icosahedron, which provides a reasonable approximation of the sphere. By choosing to work with this very regular manifold, we are able to implement the gauge equivariant convolution using a single conv2d call, making it a highly scalable and practical alternative to Spherical CNNs. Using this method, we demonstrate substantial improvements over previous methods on the task of segmenting omnidirectional images and global climate patterns.
%"""

%References 

A gauge is an assignment of a local frame of reference to each point on a manifold. A gauge transformation is a continuous set of transformations, one for each point on a manifold. Gauge equivariant convolutional networks\cite{cohen2019gauge} are a form of equivariant network that satisfy a local equivariance to the gauge element at each point on the manifold. Gauge equivariant networks have since found use cases in physical systems with natural choice of gauge.

\section{Gauge Transformations}

Let $M$ be a manifold. A gauge  is given by a position dependent map $w_p$.
\begin{align}
    w_p : \mathbb{R}^d \to T_p M
\end{align}
This map determines a reference frame is given by mapping the basis set $e_1,\dotsc, e_d$ through the gauge
\begin{align}
    w_p(e_1),\dotsc, w_p(e_d)
\end{align}
A local gauge transformation is given by a position dependent linear transformation.
\begin{align}
    g_p \in GL(d, \mathbb{R})
\end{align}
The change of frame/basis transformation is given by
\begin{align}
    w_p \mapsto w_p g_p
\end{align}
The gauge transformation provides a method for changing the reference frame at each point. The new basis after a transformation is given by
\begin{align}
    w_p(e_1)g_p,\dotsc, w_p(e_d)g_p
\end{align}
Suppose we have some component vector $v \in \mathbb{R}^d$. Under a change of basis, this vector transforms
\begin{align}
    v \mapsto g_p^{-1} v
\end{align}
so that the composition
\begin{align}
(w_p g_p)(g_p^{-1} v) &= w_p v
\end{align}
remains unchanged. The group $G = GL(d, \mathbb{R})$ is referred to as the structure group for the theory. In practice, we often choose to restrict $G$ to $SO(d)$ so that the gauge transformation is orthonormal.
%\textcolor{red}{TODO: This doesn't quite make sense? I'm getting the terms on opposite sides and they don't cancel. }

The exponential map provides transport between different points on the manifold.
\begin{align}
    \textrm{exp}_p: T_p M \to M
\end{align}
Conceptually, this map takes a point in the tangent space $T_p M$ and follows it for one (infinitesimal) unit of distance to arrive at a new point on the manifold.

\section{Gauge Equivariant Convolution for Scalar Fields}

We form an idealized mathematical model of a filter as a function

\begin{align}
    K: \mathbb{R}^d \to \mathbb{R}
\end{align}
The gauge $w_p$ lets us map $\mathbb{R}^d$ to $T_p M$. Let $f: M \to \mathbb{R}$ be a signal on the manifold. Then we can convolve filter $K$ against $f$ by the formula
\begin{align}
    (K * f)(p) &= \int_{\mathbb{R}^d} K(v) f(\textrm{exp}_p w_p(v)) dv
\end{align}
What happens if we change the gauge? To make $K * f$ invariant to gauge transformations, we must have that
\begin{align}
K(g^{-1} v) &= K(v)
\end{align}
If the structure group is $SO(d)$, then we are mandating that filter $K$ must be rotationally invariant.

\section{General Gauge Equivariant Convolutions}

\begin{align}
K(v): \mathbb{R}^{C_{in}} \to \mathbb{R}^{C_{out}}
\end{align}

\begin{align}
(K*f)(p) &= \int_{\mathbb{R}^d} K(v) \rho_{in}(g_{p \leftarrow \exp v}) f(\exp v) dv
\end{align}

The gauge is invariant if and only if
\begin{align}
\forall g \in G, \  K(g^{-1}v) &= \rho_{out}(g^{-1})K(v)\rho_{in}(g)
\end{align}

\section{Efficiency Considerations}
In the case where the gauge is chosen to be the icosahedron, we can consider the gauge equivariant convolution as an approximation to a spherical convolution.

\end{document}